\section{从一个 token 到“长出来”的内部程序}

上一章说明了为什么只看输出不够；本章先固定研究对象：一次 token 生成到底发生什么。只有把自回归概率、隐藏状态、注意力与 MLP 的角色分开，后面的 feature、transcoder 和 attribution graph 才不会沦为模糊比喻。

\subsection{自回归生成是重复前向传播}

给定输入 token 序列 $x_{1:n}$，模型按自回归分解生成后续序列：
\[
p_\theta(x_{n+1:n+T}\mid x_{1:n})
=\prod_{t=1}^{T}p_\theta(x_{n+t}\mid x_{1:n+t-1}).
\]
其中 $\theta$ 是训练后固定的参数，$T$ 是新生成 token 数，$x_{1:n+t-1}$ 是第 $t$ 次预测可读取的完整前缀。每增加一个 token，都要重新运行一次前向传播；“一次回答”因此包含许多离散计算步。

\lecturefigure{slide-05-how-one-token-is-produced.jpg}{问题入口：一个大语言模型怎样生成一个 token}{00:08:22}

\lecturefigure{slide-06-autoregressive-forward-passes.jpg}{每生成一个新 token，都以更长前缀再次运行模型}{00:08:50}

图 5 提出粒度问题，图 6 则把回答拆成多次 pass。这里要区分两种并行：序列维度上的下一个 token 仍按时间推进；但每一次 pass 内部，成千上万的 feature 与矩阵运算可同时发生。本讲所谓“parallel computation”主要指后一种，而不是一次性生成完整句子。

\begin{lstlisting}[language=Python,caption={自回归解码的最小伪代码：每个 token 对应一次新的前向传播}]
context = tokenize(prompt)
for step in range(max_new_tokens):
    logits, internal_state = transformer(context)
    next_token = sample(softmax(logits[-1]))
    context.append(next_token)
\end{lstlisting}

\subsection{隐藏状态把前缀压成下一步分布}

单次前向传播中，第 $\ell$ 层残差流可写成一个简化更新：
\[
h_i^{(\ell+1)}=h_i^{(\ell)}
+\operatorname{Attn}^{(\ell)}_i(H^{(\ell)})
+\operatorname{MLP}^{(\ell)}(h_i^{(\ell)}),
\qquad
p(x_{n+1}=v)=\operatorname{softmax}(W_U h_n^{(L)})_v.
\]
这里 $h_i^{(\ell)}$ 是位置 $i$ 在第 $\ell$ 层的 residual-stream（残差流）向量，$H^{(\ell)}$ 是所有位置的向量集合，$\operatorname{Attn}$ 负责跨位置选择与搬运信息，$\operatorname{MLP}$ 负责逐位置非线性变换，$W_U$ 把最后一层状态映射为词表 logits，$v$ 是候选 token。

\lecturefigure{slide-07-single-forward-pass-vectors.jpg}{输入向量经过隐藏层，最终形成下一 token 的分数分布}{00:09:42}

读图时先看左侧 token 如何被编码为向量，再看中间层怎样反复重写表示，最后看右侧每个候选词的 score。输出概率只是最终投影；真正需要解释的是中间层如何让 “assist” 获得更高分。一个 attribution graph 就是在尝试恢复这条从输入 feature 到输出 logit 的影响链。

\begin{knowledgebox}{术语第一次出现：residual stream、neuron 与 feature}
\textbf{Residual stream} 是各层共享的高维状态通道；attention 与 MLP 都向其中读写。\textbf{Neuron} 通常指 MLP 隐层中的单个坐标。\textbf{Feature} 则是研究者试图识别的、在激活空间中代表某种概念或操作的方向；一个 feature 可以跨越多个 neuron，一个 neuron 也可能参与多个 feature。
\end{knowledgebox}

\subsection{“长出来”不是神秘主义，而是工程事实}

架构规定了可学习计算的容器，训练数据与梯度下降决定容器中最终形成什么算法。工程师知道矩阵维度、层数和损失函数，却通常不知道模型为何在某个 prompt 上激活 Dallas、Texas、capital 与 Austin 的特定组合。逆向分析的目标，是把这种“我们写了训练过程，但没有写内部程序”的缺口变成可检验对象。

\lecturefigure{slide-08-models-grown-not-built.jpg}{核心比喻：语言模型是 grown，而不是逐模块 built}{00:09:51}

\lecturefigure{slide-09-architecture-versus-grown-model.jpg}{人类搭建架构，训练过程让内部机制以难以直接阅读的方式生长}{00:10:20}

两图的证据边界很重要：它们不意味着模型不可理解，也不意味着任何拟人化解释都成立。恰恰相反，“grown”要求更严格的 reverse engineering：提出部件、测量连接、做干预、记录失败，并承认当前方法只能解释一部分计算。

\begin{importantbox}{“Grown not built”的可操作含义}
我们可以精确知道训练代码，却仍不知道训练后网络在单个输入上采用哪套内部算法。可解释性因此关注\textbf{post-training program discovery}：在固定权重上恢复局部、输入条件化的计算结构。
\end{importantbox}

\teachervoice{00:09:47--00:10:20，老师用 grown, not built 区分架构设计与内部算法形成。课堂提示：这不是把模型神秘化，而是说明传统代码审查无法直接告诉我们训练后网络在做什么。}

\subsection{本章小结}

一次回答由多次 token 级前向传播构成，而一次前向传播内部又包含大规模并行计算。后续方法需要解释的，不是 softmax 最终选了哪个词，而是隐藏状态中哪些 feature 被激活、它们怎样通过 attention 与 MLP 共同改变输出。

